\documentclass[a4paper]{article}
\usepackage[utf8]{inputenc} 	% codificacao de caracteres
\usepackage[T1]{fontenc}    	% codificacao de fontes
\usepackage[brazil]{babel}  	% idioma
\usepackage[table]{xcolor}    	% linhas coloridas alternadamente
\usepackage{graphicx}
\usepackage{multirow}
\usepackage{soulutf8}			% cor de fundo do texto que aceita utf8
\usepackage{makecell}			% mais de uma linha na célula da tabela

% espaçamento paragrafos
\setlength{\parindent}{0em}
%\setlength{\parskip}{2em}
%\renewcommand{\baselinestretch}{1}

%opening
\title{Universidade Federal do Ceará - Crateús}
\author{Ítalo Ribeiro}
\date{}

\begin{document}

\fontfamily{cmss}
\fontsize{12}{18}
\selectfont

\begin{center}
	\includegraphics[scale=0.18]{brasao2_vertical}
\end{center}

\begin{center}
	\vspace{14ex}
	Campus de Crateús\\	
	Estágio Supervisionado\\
	\vspace{10ex}
	\textbf{{
			\large PLANO DE TRABALHO DE ESTÁGIO SUPERVISIONADO
	}}\\
	\vspace{16ex}
	{\large 
		Atividade de Estágio Supervisionado\\
		Estágio <Especificar Instituição>\\
		<Título do Projeto que está sendo desenvolvido>
	}
\end{center}

\newpage

\section*{Identificação}

\textbf{Estagiário:}\\

Nome:  \\
Matrícula:  \\
Telefones:  \\
E-mail:  \\
Cargo:  \\

\textbf{Instituição do Estágio:}\\

Nome:  \\
Endereço:  \\
Ramo de Atividade:  \\
Supervisor:  \\
Telefone:  \\
E-mail:  \\

\textbf{Professor Orientador:}

\newpage

\tableofcontents

\newpage

\addcontentsline{toc}{section}{Introdução}
\section*{Introdução}

\sethlcolor{lightgray}

Este documento tem por objetivo mostrar o trabalho desenvolvido pelo aluno(a) \hl{<Nome do aluno(a)>} durante o estágio no(a) \hl{<Nome da empresa>}.

\addcontentsline{toc}{section}{Empresa}
\section*{Empresa}

\hl{<Adicionar suas impressões em relação à empresa. Pontos positivos e negativos do local de estágio> OBS.: Não é aquele texto padrão da empresa, são suas impressões. \\
<Relato sobre processos e ferramentas adotados pelas empresas.>
}

\addcontentsline{toc}{section}{Relato de Experiência}
\section*{Relato de Experiência}

\hl{<Fazer um relato de experiência descrevendo aqui tudo como foi seu estágio.>}

\subsection*{Cronograma de atividades}

\hl{<As atividades descritas no relato de experiência acima devem ser demonstradas aqui em uma tabela com o cronograma destas atividades.>}

\subsection*{Lições Aprendidas}

\hl{<Citar as lições aprendidas durante o desenvolvimento das atividades e justificar a importância destas lições aprendidas.>}

\subsection*{Dificuldades Encontradas}

\hl{<Citar as dificuldades encontradas durante o desenvolvimento das atividades e mostrar se elas foram superadas e como foram superadas.>}

\subsection*{Oportunidades de Melhoria}

\hl{<Nada está tão bom que não possa melhorar. Pense como você pode melhorar, o que pode aprender para desempenhar suas atividades de forma mais satisfatória ainda.>}


\addcontentsline{toc}{section}{Conclusão}
\section*{Conclusão}

\hl{<Faça um encerramento de seu relatório.>}


\end{document}
